\documentclass[11pt,a4paper]{article}

% Preambulo compartido por todas las partes del informe.
% Cada parte hace % Preambulo compartido por todas las partes del informe.
% Cada parte hace % Preambulo compartido por todas las partes del informe.
% Cada parte hace \input{preambulo} y despues carga sus propias tablas.

\usepackage[utf8]{inputenc}
\usepackage[T1]{fontenc}
\usepackage[spanish,es-noquoting]{babel}
\usepackage[a4paper,margin=20mm]{geometry}
\usepackage{amsmath,amssymb}
\usepackage{booktabs}
\usepackage{array}
\usepackage{xcolor}
\usepackage{enumitem}
\usepackage{titlesec}

% ------------------------------------------------ colores
% 'gris' lo usan las tablas generadas para las entradas repetidas por simetria
\definecolor{gris}{RGB}{170,170,170}

% ------------------------------------------------ espaciado
\titleformat{\section}{\normalfont\large\bfseries}{\thesection}{0.6em}{}
\titlespacing*{\section}{0pt}{2.4ex plus 1ex}{1.2ex}
\titleformat{\subsection}{\normalfont\bfseries}{\thesubsection}{0.6em}{}
\titlespacing*{\subsection}{0pt}{1.8ex plus .5ex}{0.8ex}

\setlength{\parindent}{0pt}
\setlength{\parskip}{0.6ex}
\renewcommand{\arraystretch}{1.08}

% ------------------------------------------------ atajos
\newcommand{\hextt}[1]{\texttt{#1}}
 y despues carga sus propias tablas.

\usepackage[utf8]{inputenc}
\usepackage[T1]{fontenc}
\usepackage[spanish,es-noquoting]{babel}
\usepackage[a4paper,margin=20mm]{geometry}
\usepackage{amsmath,amssymb}
\usepackage{booktabs}
\usepackage{array}
\usepackage{xcolor}
\usepackage{enumitem}
\usepackage{titlesec}

% ------------------------------------------------ colores
% 'gris' lo usan las tablas generadas para las entradas repetidas por simetria
\definecolor{gris}{RGB}{170,170,170}

% ------------------------------------------------ espaciado
\titleformat{\section}{\normalfont\large\bfseries}{\thesection}{0.6em}{}
\titlespacing*{\section}{0pt}{2.4ex plus 1ex}{1.2ex}
\titleformat{\subsection}{\normalfont\bfseries}{\thesubsection}{0.6em}{}
\titlespacing*{\subsection}{0pt}{1.8ex plus .5ex}{0.8ex}

\setlength{\parindent}{0pt}
\setlength{\parskip}{0.6ex}
\renewcommand{\arraystretch}{1.08}

% ------------------------------------------------ atajos
\newcommand{\hextt}[1]{\texttt{#1}}
 y despues carga sus propias tablas.

\usepackage[utf8]{inputenc}
\usepackage[T1]{fontenc}
\usepackage[spanish,es-noquoting]{babel}
\usepackage[a4paper,margin=20mm]{geometry}
\usepackage{amsmath,amssymb}
\usepackage{booktabs}
\usepackage{array}
\usepackage{xcolor}
\usepackage{enumitem}
\usepackage{titlesec}

% ------------------------------------------------ colores
% 'gris' lo usan las tablas generadas para las entradas repetidas por simetria
\definecolor{gris}{RGB}{170,170,170}

% ------------------------------------------------ espaciado
\titleformat{\section}{\normalfont\large\bfseries}{\thesection}{0.6em}{}
\titlespacing*{\section}{0pt}{2.4ex plus 1ex}{1.2ex}
\titleformat{\subsection}{\normalfont\bfseries}{\thesubsection}{0.6em}{}
\titlespacing*{\subsection}{0pt}{1.8ex plus .5ex}{0.8ex}

\setlength{\parindent}{0pt}
\setlength{\parskip}{0.6ex}
\renewcommand{\arraystretch}{1.08}

% ------------------------------------------------ atajos
\newcommand{\hextt}[1]{\texttt{#1}}


% generado por gen_parte_a.py -- no editar a mano
\newcommand{\parEj}{\texttt{9A5}}
\newcommand{\cwEj}{\texttt{A5C9A5}}

% generado por gen_parte_a.py -- no editar a mano
\newcommand{\rAXor}{%
\begin{tabular}{@{}r@{\;\;}l@{\;\;}l@{}}
 & $b_{0}$ & \texttt{1001\,1000\,1111} \\
$\oplus$ & $b_{2}$ & \texttt{0011\,0101\,0111} \\
\cmidrule(lr){2-3}
 & & \texttt{1010\,1101\,1000} \\
$\oplus$ & $b_{5}$ & \texttt{0111\,1100\,1100} \\
\cmidrule(lr){2-3}
 & & \texttt{1101\,0001\,0100} \\
$\oplus$ & $b_{7}$ & \texttt{0010\,1011\,1110} \\
\cmidrule(lr){2-3}
 & & \texttt{1111\,1010\,1010} \\
$\oplus$ & $b_{8}$ & \texttt{1000\,0111\,1011} \\
\cmidrule(lr){2-3}
 & & \texttt{0111\,1101\,0001} \\
$\oplus$ & $b_{9}$ & \texttt{1110\,0111\,0100} \\
\cmidrule(lr){2-3}
 & & \texttt{1001\,1010\,0101} \\
$\oplus$ & $b_{11}$ & \texttt{1110\,1010\,1001} \\
\cmidrule(lr){2-3}
 & $B\,r_{[23:12]}$ & \texttt{0111\,0000\,1100} \;$=$\; \texttt{70C} \\
\end{tabular}
}
\newcommand{\rASind}{%
\begin{tabular}{@{}r@{\;\;}l@{\;\;}l@{}}
 & $B\,r_{[23:12]}$ & \texttt{0111\,0000\,1100} \\
$\oplus$ & $r_{[11:0]}$ & \texttt{1001\,1010\,0110} \\
\midrule
 & $s$ & \texttt{1110\,1010\,1010} \;$=$\; \texttt{EAA} \\
\end{tabular}
}
\newcommand{\rAPesoS}{7}
\newcommand{\rATablaDos}{%
\begin{tabular}{@{}cccc@{}}
\toprule
$i$ & $b_i$ & $s\oplus b_i$ & $w$ \\
\midrule
0 & \texttt{1001\,1000\,1111} & \texttt{0111\,0010\,0101} & 6 \\
1 & \texttt{0100\,1110\,0111} & \texttt{1010\,0100\,1101} & 6 \\
2 & \texttt{0011\,0101\,0111} & \texttt{1101\,1111\,1101} & 10 \\
3 & \texttt{1011\,1110\,0010} & \texttt{0101\,0100\,1000} & 4 \\
4 & \texttt{1101\,1101\,0001} & \texttt{0011\,0111\,1011} & 8 \\
5 & \texttt{0111\,1100\,1100} & \texttt{1001\,0110\,0110} & 6 \\
6 & \texttt{0101\,0011\,1101} & \texttt{1011\,1001\,0111} & 8 \\
7 & \texttt{0010\,1011\,1110} & \texttt{1100\,0001\,0100} & 4 \\
8 & \texttt{1000\,0111\,1011} & \texttt{0110\,1101\,0001} & 6 \\
9 & \texttt{1110\,0111\,0100} & \texttt{0000\,1101\,1110} & 6 \\
10 & \texttt{1111\,0001\,1010} & \texttt{0001\,1011\,0000} & 4 \\
11 & \texttt{1110\,1010\,1001} & \texttt{0000\,0000\,0011} & $\mathbf{2}$ $\leftarrow$ \\
\bottomrule
\end{tabular}
}
\newcommand{\rARes}{%
\begin{tabular}{@{}r@{\;\;}l@{\;\;}l@{}}
 & $r$ & \texttt{1010\,0101\,1101\,1001\,1010\,0110} \\
$\oplus$ & $e$ & \texttt{0000\,0000\,0001\,0000\,0000\,0011} \\
\midrule
 & $v$ & \texttt{1010\,0101\,1100\,1001\,1010\,0101} \;$=$\; \texttt{A5C9A5} \\
\end{tabular}
}

% generado por gen_parte_a.py -- no editar a mano
\newcommand{\rBXor}{%
\begin{tabular}{@{}r@{\;\;}l@{\;\;}l@{}}
 & $b_{0}$ & \texttt{1001\,1000\,1111} \\
$\oplus$ & $b_{2}$ & \texttt{0011\,0101\,0111} \\
\cmidrule(lr){2-3}
 & & \texttt{1010\,1101\,1000} \\
$\oplus$ & $b_{5}$ & \texttt{0111\,1100\,1100} \\
\cmidrule(lr){2-3}
 & & \texttt{1101\,0001\,0100} \\
$\oplus$ & $b_{7}$ & \texttt{0010\,1011\,1110} \\
\cmidrule(lr){2-3}
 & & \texttt{1111\,1010\,1010} \\
$\oplus$ & $b_{8}$ & \texttt{1000\,0111\,1011} \\
\cmidrule(lr){2-3}
 & & \texttt{0111\,1101\,0001} \\
$\oplus$ & $b_{9}$ & \texttt{1110\,0111\,0100} \\
\cmidrule(lr){2-3}
 & & \texttt{1001\,1010\,0101} \\
$\oplus$ & $b_{10}$ & \texttt{1111\,0001\,1010} \\
\cmidrule(lr){2-3}
 & & \texttt{0110\,1011\,1111} \\
$\oplus$ & $b_{11}$ & \texttt{1110\,1010\,1001} \\
\cmidrule(lr){2-3}
 & $B\,r_{[23:12]}$ & \texttt{1000\,0001\,0110} \;$=$\; \texttt{816} \\
\end{tabular}
}
\newcommand{\rBSind}{%
\begin{tabular}{@{}r@{\;\;}l@{\;\;}l@{}}
 & $B\,r_{[23:12]}$ & \texttt{1000\,0001\,0110} \\
$\oplus$ & $r_{[11:0]}$ & \texttt{1001\,1010\,0100} \\
\midrule
 & $s$ & \texttt{0001\,1011\,0010} \;$=$\; \texttt{1B2} \\
\end{tabular}
}
\newcommand{\rBPesoS}{5}
\newcommand{\rBTablaDos}{%
\begin{tabular}{@{}cccc@{}}
\toprule
$i$ & $b_i$ & $s\oplus b_i$ & $w$ \\
\midrule
0 & \texttt{1001\,1000\,1111} & \texttt{1000\,0011\,1101} & 6 \\
1 & \texttt{0100\,1110\,0111} & \texttt{0101\,0101\,0101} & 6 \\
2 & \texttt{0011\,0101\,0111} & \texttt{0010\,1110\,0101} & 6 \\
3 & \texttt{1011\,1110\,0010} & \texttt{1010\,0101\,0000} & 4 \\
4 & \texttt{1101\,1101\,0001} & \texttt{1100\,0110\,0011} & 6 \\
5 & \texttt{0111\,1100\,1100} & \texttt{0110\,0111\,1110} & 8 \\
6 & \texttt{0101\,0011\,1101} & \texttt{0100\,1000\,1111} & 6 \\
7 & \texttt{0010\,1011\,1110} & \texttt{0011\,0000\,1100} & 4 \\
8 & \texttt{1000\,0111\,1011} & \texttt{1001\,1100\,1001} & 6 \\
9 & \texttt{1110\,0111\,0100} & \texttt{1111\,1100\,0110} & 8 \\
10 & \texttt{1111\,0001\,1010} & \texttt{1110\,1010\,1000} & 6 \\
11 & \texttt{1110\,1010\,1001} & \texttt{1111\,0001\,1011} & 8 \\
\bottomrule
\end{tabular}
}
\newcommand{\rBQ}{%
\begin{tabular}{@{}r@{\;\;}l@{\;\;}l@{}}
 & $b_{3}$ & \texttt{1011\,1110\,0010} \\
$\oplus$ & $b_{4}$ & \texttt{1101\,1101\,0001} \\
\cmidrule(lr){2-3}
 & & \texttt{0110\,0011\,0011} \\
$\oplus$ & $b_{6}$ & \texttt{0101\,0011\,1101} \\
\cmidrule(lr){2-3}
 & & \texttt{0011\,0000\,1110} \\
$\oplus$ & $b_{7}$ & \texttt{0010\,1011\,1110} \\
\cmidrule(lr){2-3}
 & & \texttt{0001\,1011\,0000} \\
$\oplus$ & $b_{10}$ & \texttt{1111\,0001\,1010} \\
\cmidrule(lr){2-3}
 & $q$ & \texttt{1110\,1010\,1010} \;$=$\; \texttt{EAA} \\
\end{tabular}
}
\newcommand{\rBPesoQ}{7}
\newcommand{\rBTablaCuatro}{%
\begin{tabular}{@{}cccc@{}}
\toprule
$i$ & $b_i$ & $q\oplus b_i$ & $w$ \\
\midrule
0 & \texttt{1001\,1000\,1111} & \texttt{0111\,0010\,0101} & 6 \\
1 & \texttt{0100\,1110\,0111} & \texttt{1010\,0100\,1101} & 6 \\
2 & \texttt{0011\,0101\,0111} & \texttt{1101\,1111\,1101} & 10 \\
3 & \texttt{1011\,1110\,0010} & \texttt{0101\,0100\,1000} & 4 \\
4 & \texttt{1101\,1101\,0001} & \texttt{0011\,0111\,1011} & 8 \\
5 & \texttt{0111\,1100\,1100} & \texttt{1001\,0110\,0110} & 6 \\
6 & \texttt{0101\,0011\,1101} & \texttt{1011\,1001\,0111} & 8 \\
7 & \texttt{0010\,1011\,1110} & \texttt{1100\,0001\,0100} & 4 \\
8 & \texttt{1000\,0111\,1011} & \texttt{0110\,1101\,0001} & 6 \\
9 & \texttt{1110\,0111\,0100} & \texttt{0000\,1101\,1110} & 6 \\
10 & \texttt{1111\,0001\,1010} & \texttt{0001\,1011\,0000} & 4 \\
11 & \texttt{1110\,1010\,1001} & \texttt{0000\,0000\,0011} & $\mathbf{2}$ $\leftarrow$ \\
\bottomrule
\end{tabular}
}
\newcommand{\rBRes}{%
\begin{tabular}{@{}r@{\;\;}l@{\;\;}l@{}}
 & $r$ & \texttt{1010\,0101\,1111\,1001\,1010\,0100} \\
$\oplus$ & $e$ & \texttt{0000\,0000\,0011\,0000\,0000\,0001} \\
\midrule
 & $v$ & \texttt{1010\,0101\,1100\,1001\,1010\,0101} \;$=$\; \texttt{A5C9A5} \\
\end{tabular}
}

% generado por gen_parte_a.py -- no editar a mano
\newcommand{\rCXor}{%
\begin{tabular}{@{}r@{\;\;}l@{\;\;}l@{}}
 & $b_{0}$ & \texttt{1001\,1000\,1111} \\
$\oplus$ & $b_{2}$ & \texttt{0011\,0101\,0111} \\
\cmidrule(lr){2-3}
 & & \texttt{1010\,1101\,1000} \\
$\oplus$ & $b_{5}$ & \texttt{0111\,1100\,1100} \\
\cmidrule(lr){2-3}
 & & \texttt{1101\,0001\,0100} \\
$\oplus$ & $b_{7}$ & \texttt{0010\,1011\,1110} \\
\cmidrule(lr){2-3}
 & & \texttt{1111\,1010\,1010} \\
$\oplus$ & $b_{8}$ & \texttt{1000\,0111\,1011} \\
\cmidrule(lr){2-3}
 & & \texttt{0111\,1101\,0001} \\
$\oplus$ & $b_{9}$ & \texttt{1110\,0111\,0100} \\
\cmidrule(lr){2-3}
 & $B\,r_{[23:12]}$ & \texttt{1001\,1010\,0101} \;$=$\; \texttt{9A5} \\
\end{tabular}
}
\newcommand{\rCSind}{%
\begin{tabular}{@{}r@{\;\;}l@{\;\;}l@{}}
 & $B\,r_{[23:12]}$ & \texttt{1001\,1010\,0101} \\
$\oplus$ & $r_{[11:0]}$ & \texttt{1001\,1010\,1010} \\
\midrule
 & $s$ & \texttt{0000\,0000\,1111} \;$=$\; \texttt{00F} \\
\end{tabular}
}
\newcommand{\rCPesoS}{4}
\newcommand{\rCTablaDos}{%
\begin{tabular}{@{}cccc@{}}
\toprule
$i$ & $b_i$ & $s\oplus b_i$ & $w$ \\
\midrule
0 & \texttt{1001\,1000\,1111} & \texttt{1001\,1000\,0000} & 3 \\
1 & \texttt{0100\,1110\,0111} & \texttt{0100\,1110\,1000} & 5 \\
2 & \texttt{0011\,0101\,0111} & \texttt{0011\,0101\,1000} & 5 \\
3 & \texttt{1011\,1110\,0010} & \texttt{1011\,1110\,1101} & 9 \\
4 & \texttt{1101\,1101\,0001} & \texttt{1101\,1101\,1110} & 9 \\
5 & \texttt{0111\,1100\,1100} & \texttt{0111\,1100\,0011} & 7 \\
6 & \texttt{0101\,0011\,1101} & \texttt{0101\,0011\,0010} & 5 \\
7 & \texttt{0010\,1011\,1110} & \texttt{0010\,1011\,0001} & 5 \\
8 & \texttt{1000\,0111\,1011} & \texttt{1000\,0111\,0100} & 5 \\
9 & \texttt{1110\,0111\,0100} & \texttt{1110\,0111\,1011} & 9 \\
10 & \texttt{1111\,0001\,1010} & \texttt{1111\,0001\,0101} & 7 \\
11 & \texttt{1110\,1010\,1001} & \texttt{1110\,1010\,0110} & 7 \\
\bottomrule
\end{tabular}
}
\newcommand{\rCQ}{%
\begin{tabular}{@{}r@{\;\;}l@{\;\;}l@{}}
 & $b_{8}$ & \texttt{1000\,0111\,1011} \\
$\oplus$ & $b_{9}$ & \texttt{1110\,0111\,0100} \\
\cmidrule(lr){2-3}
 & & \texttt{0110\,0000\,1111} \\
$\oplus$ & $b_{10}$ & \texttt{1111\,0001\,1010} \\
\cmidrule(lr){2-3}
 & & \texttt{1001\,0001\,0101} \\
$\oplus$ & $b_{11}$ & \texttt{1110\,1010\,1001} \\
\cmidrule(lr){2-3}
 & $q$ & \texttt{0111\,1011\,1100} \;$=$\; \texttt{7BC} \\
\end{tabular}
}
\newcommand{\rCPesoQ}{8}
\newcommand{\rCTablaCuatro}{%
\begin{tabular}{@{}cccc@{}}
\toprule
$i$ & $b_i$ & $q\oplus b_i$ & $w$ \\
\midrule
0 & \texttt{1001\,1000\,1111} & \texttt{1110\,0011\,0011} & 7 \\
1 & \texttt{0100\,1110\,0111} & \texttt{0011\,0101\,1011} & 7 \\
2 & \texttt{0011\,0101\,0111} & \texttt{0100\,1110\,1011} & 7 \\
3 & \texttt{1011\,1110\,0010} & \texttt{1100\,0101\,1110} & 7 \\
4 & \texttt{1101\,1101\,0001} & \texttt{1010\,0110\,1101} & 7 \\
5 & \texttt{0111\,1100\,1100} & \texttt{0000\,0111\,0000} & 3 \\
6 & \texttt{0101\,0011\,1101} & \texttt{0010\,1000\,0001} & 3 \\
7 & \texttt{0010\,1011\,1110} & \texttt{0101\,0000\,0010} & 3 \\
8 & \texttt{1000\,0111\,1011} & \texttt{1111\,1100\,0111} & 9 \\
9 & \texttt{1110\,0111\,0100} & \texttt{1001\,1100\,1000} & 5 \\
10 & \texttt{1111\,0001\,1010} & \texttt{1000\,1010\,0110} & 5 \\
11 & \texttt{1110\,1010\,1001} & \texttt{1001\,0001\,0101} & 5 \\
\bottomrule
\end{tabular}
}
\newcommand{\rCRes}{%
---
}



\begin{document}

\begin{center}
  {\large\bfseries Código de Golay extendido $(24,12)$ --- Parte A}\\[0.3ex]
  {\small Trabajo Práctico 2 · Curso de FEC 2026 · Fundación Fulgor}\\[0.3ex]
  {\small Villar, Pedro}
\end{center}


\section*{Datos}

Matriz $B$ (simétrica, $B^2=I$), idéntica para todos los grupos:

\begin{center}\small
% generado por gen_parte_a.py -- no editar a mano
\begin{tabular}{@{}llc@{\hspace{2.5em}}llc@{}}
\toprule
fila & hex & binario & fila & hex & binario \\
\midrule
$b_{0}$ & \texttt{98F} & \texttt{1001\,1000\,1111} & $b_{6}$ & \texttt{53D} & \texttt{0101\,0011\,1101} \\
$b_{1}$ & \texttt{4E7} & \texttt{0100\,1110\,0111} & $b_{7}$ & \texttt{2BE} & \texttt{0010\,1011\,1110} \\
$b_{2}$ & \texttt{357} & \texttt{0011\,0101\,0111} & $b_{8}$ & \texttt{87B} & \texttt{1000\,0111\,1011} \\
$b_{3}$ & \texttt{BE2} & \texttt{1011\,1110\,0010} & $b_{9}$ & \texttt{E74} & \texttt{1110\,0111\,0100} \\
$b_{4}$ & \texttt{DD1} & \texttt{1101\,1101\,0001} & $b_{10}$ & \texttt{F1A} & \texttt{1111\,0001\,1010} \\
$b_{5}$ & \texttt{7CC} & \texttt{0111\,1100\,1100} & $b_{11}$ & \texttt{EA9} & \texttt{1110\,1010\,1001} \\
\bottomrule
\end{tabular}

\end{center}

Convención de bits: $cw[23{:}12]=$ mensaje, $cw[11{:}0]=$ paridad; el bit
$m[11-i]$ es $m_i$ y el bit $[11-j]$ de $b_i$ es $B[i][j]$. Con esa
convención, multiplicar un vector $x$ de 12 bits por $B$ es sumar
(XOR) las filas $b_i$ tales que $x_i=1$:
\[
  x\cdot B \;=\; \bigoplus_{i\,:\,x_i=1} b_i .
\]

\section{Ejercicio 1 --- Matrices y parámetros}

\subsection{$G$ y $H$}

$G=[\,I_{12}\mid B\,]$ y $H=[\,B\mid I_{12}\,]$, ambas $12\times 24$:

\begin{center}\small
\begin{minipage}[t]{0.47\textwidth}
  \centering $G$\\[0.6ex]
  % generado por gen_parte_a.py -- no editar a mano
\begin{tabular}{@{}r@{\;\;}c@{\;\;}c@{}}
$g_{0}$ & \texttt{100000000000} & \texttt{100110001111} \\
$g_{1}$ & \texttt{010000000000} & \texttt{010011100111} \\
$g_{2}$ & \texttt{001000000000} & \texttt{001101010111} \\
$g_{3}$ & \texttt{000100000000} & \texttt{101111100010} \\
$g_{4}$ & \texttt{000010000000} & \texttt{110111010001} \\
$g_{5}$ & \texttt{000001000000} & \texttt{011111001100} \\
$g_{6}$ & \texttt{000000100000} & \texttt{010100111101} \\
$g_{7}$ & \texttt{000000010000} & \texttt{001010111110} \\
$g_{8}$ & \texttt{000000001000} & \texttt{100001111011} \\
$g_{9}$ & \texttt{000000000100} & \texttt{111001110100} \\
$g_{10}$ & \texttt{000000000010} & \texttt{111100011010} \\
$g_{11}$ & \texttt{000000000001} & \texttt{111010101001} \\
\end{tabular}

\end{minipage}\hfill
\begin{minipage}[t]{0.47\textwidth}
  \centering $H$\\[0.6ex]
  % generado por gen_parte_a.py -- no editar a mano
\begin{tabular}{@{}r@{\;\;}c@{\;\;}c@{}}
$h_{0}$ & \texttt{100110001111} & \texttt{100000000000} \\
$h_{1}$ & \texttt{010011100111} & \texttt{010000000000} \\
$h_{2}$ & \texttt{001101010111} & \texttt{001000000000} \\
$h_{3}$ & \texttt{101111100010} & \texttt{000100000000} \\
$h_{4}$ & \texttt{110111010001} & \texttt{000010000000} \\
$h_{5}$ & \texttt{011111001100} & \texttt{000001000000} \\
$h_{6}$ & \texttt{010100111101} & \texttt{000000100000} \\
$h_{7}$ & \texttt{001010111110} & \texttt{000000010000} \\
$h_{8}$ & \texttt{100001111011} & \texttt{000000001000} \\
$h_{9}$ & \texttt{111001110100} & \texttt{000000000100} \\
$h_{10}$ & \texttt{111100011010} & \texttt{000000000010} \\
$h_{11}$ & \texttt{111010101001} & \texttt{000000000001} \\
\end{tabular}

\end{minipage}
\end{center}

\subsection{$B$ es simétrica}

Se lee cada columna $j$ de $B$ de arriba a abajo y se compara con la fila
$b_j$. El bit $i$ de la columna $j$ es $B[i][j]$, o sea el bit $[11-j]$ de
$b_i$:

\begin{center}\small
% generado por gen_parte_a.py -- no editar a mano
\begin{tabular}{@{}ccc@{\hspace{2.5em}}ccc@{}}
\toprule
$j$ & columna $j$ & $b_j$ & $j$ & columna $j$ & $b_j$ \\
\midrule
0 & \texttt{1001\,1000\,1111} & \texttt{1001\,1000\,1111} & 6 & \texttt{0101\,0011\,1101} & \texttt{0101\,0011\,1101} \\
1 & \texttt{0100\,1110\,0111} & \texttt{0100\,1110\,0111} & 7 & \texttt{0010\,1011\,1110} & \texttt{0010\,1011\,1110} \\
2 & \texttt{0011\,0101\,0111} & \texttt{0011\,0101\,0111} & 8 & \texttt{1000\,0111\,1011} & \texttt{1000\,0111\,1011} \\
3 & \texttt{1011\,1110\,0010} & \texttt{1011\,1110\,0010} & 9 & \texttt{1110\,0111\,0100} & \texttt{1110\,0111\,0100} \\
4 & \texttt{1101\,1101\,0001} & \texttt{1101\,1101\,0001} & 10 & \texttt{1111\,0001\,1010} & \texttt{1111\,0001\,1010} \\
5 & \texttt{0111\,1100\,1100} & \texttt{0111\,1100\,1100} & 11 & \texttt{1110\,1010\,1001} & \texttt{1110\,1010\,1001} \\
\bottomrule
\end{tabular}

\end{center}

Las doce coinciden, entonces $B^{T}=B$.

\subsection{$G\,H^{T}=0$}

Con $H=[\,B\mid I_{12}\,]$ se tiene $H^{T}=\begin{bmatrix}B^{T}\\ I_{12}\end{bmatrix}$
(bloques de $12\times 12$), y el producto se parte en dos sumandos:
\[
  G\,H^{T}
  \;=\;
  \begin{bmatrix} I_{12} & B \end{bmatrix}
  \begin{bmatrix} B^{T} \\ I_{12} \end{bmatrix}
  \;=\; I_{12}\,B^{T} + B\,I_{12}
  \;=\; B^{T}+B
  \;=\; B+B \;=\; 0 ,
\]
donde la anteúltima igualdad es la simetría recién verificada y la última
vale porque en $GF(2)$ la suma es XOR y $x+x=0$.

Chequeo numérico sobre una fila: $g_0=(u_0\mid b_0)$, entonces
$g_0H^{T}=u_0B^{T}+b_0=b_0+b_0=0$, y lo mismo para las otras once.

\subsection{$B^{2}=I$}

Como $B$ es simétrica, $B^2=B\,B^{T}$, así que la entrada $(i,j)$ es el
producto interno de dos \emph{filas}:
\[
  (B^2)_{ij} \;=\; \sum_{k=0}^{11} B[i][k]\,B[k][j]
  \;=\; b_i\cdot b_j
  \;=\; w\!\left(b_i \wedge b_j\right) \bmod 2 ,
\]
o sea la paridad de la cantidad de posiciones donde $b_i$ y $b_j$ tienen
ambos un 1. Además $B^2$ también es simétrica, así que alcanza con calcular
las $\binom{12}{2}+12=78$ entradas con $j\ge i$.

Tres entradas desarrolladas:

% generado por gen_parte_a.py -- no editar a mano
\begin{align*}
b_{0} &= \texttt{1001\,1000\,1111} \\
b_{0} &= \texttt{1001\,1000\,1111} \\
b_{0}\wedge b_{0} &= \texttt{1001\,1000\,1111} \quad\Longrightarrow\quad w = 7 \equiv 1 \!\pmod 2
\end{align*}
\begin{align*}
b_{0} &= \texttt{1001\,1000\,1111} \\
b_{1} &= \texttt{0100\,1110\,0111} \\
b_{0}\wedge b_{1} &= \texttt{0000\,1000\,0111} \quad\Longrightarrow\quad w = 4 \equiv 0 \!\pmod 2
\end{align*}
\begin{align*}
b_{3} &= \texttt{1011\,1110\,0010} \\
b_{7} &= \texttt{0010\,1011\,1110} \\
b_{3}\wedge b_{7} &= \texttt{0010\,1010\,0010} \quad\Longrightarrow\quad w = 4 \equiv 0 \!\pmod 2
\end{align*}


La tabla completa de $w(b_i\wedge b_j)$ (en gris, las entradas repetidas por
simetría):

\begin{center}\footnotesize
% generado por gen_parte_a.py -- no editar a mano
\begin{tabular}{@{}r|*{12}{c}@{}}
$w(b_i\wedge b_j)$ & $\mathbf{0}$ & $\mathbf{1}$ & $\mathbf{2}$ & $\mathbf{3}$ & $\mathbf{4}$ & $\mathbf{5}$ & $\mathbf{6}$ & $\mathbf{7}$ & $\mathbf{8}$ & $\mathbf{9}$ & $\mathbf{10}$ & $\mathbf{11}$ \\
\midrule
$\mathbf{0}$ & \textbf{7} & 4 & 4 & 4 & 4 & 4 & 4 & 4 & 4 & 2 & 4 & 4 \\
$\mathbf{1}$ & \textcolor{gris}{4} & \textbf{7} & 4 & 4 & 4 & 4 & 4 & 4 & 4 & 4 & 2 & 4 \\
$\mathbf{2}$ & \textcolor{gris}{4} & \textcolor{gris}{4} & \textbf{7} & 4 & 4 & 4 & 4 & 4 & 4 & 4 & 4 & 2 \\
$\mathbf{3}$ & \textcolor{gris}{4} & \textcolor{gris}{4} & \textcolor{gris}{4} & \textbf{7} & 4 & 4 & 2 & 4 & 4 & 4 & 4 & 4 \\
$\mathbf{4}$ & \textcolor{gris}{4} & \textcolor{gris}{4} & \textcolor{gris}{4} & \textcolor{gris}{4} & \textbf{7} & 4 & 4 & 2 & 4 & 4 & 4 & 4 \\
$\mathbf{5}$ & \textcolor{gris}{4} & \textcolor{gris}{4} & \textcolor{gris}{4} & \textcolor{gris}{4} & \textcolor{gris}{4} & \textbf{7} & 4 & 4 & 2 & 4 & 4 & 4 \\
$\mathbf{6}$ & \textcolor{gris}{4} & \textcolor{gris}{4} & \textcolor{gris}{4} & \textcolor{gris}{2} & \textcolor{gris}{4} & \textcolor{gris}{4} & \textbf{7} & 4 & 4 & 4 & 4 & 4 \\
$\mathbf{7}$ & \textcolor{gris}{4} & \textcolor{gris}{4} & \textcolor{gris}{4} & \textcolor{gris}{4} & \textcolor{gris}{2} & \textcolor{gris}{4} & \textcolor{gris}{4} & \textbf{7} & 4 & 4 & 4 & 4 \\
$\mathbf{8}$ & \textcolor{gris}{4} & \textcolor{gris}{4} & \textcolor{gris}{4} & \textcolor{gris}{4} & \textcolor{gris}{4} & \textcolor{gris}{2} & \textcolor{gris}{4} & \textcolor{gris}{4} & \textbf{7} & 4 & 4 & 4 \\
$\mathbf{9}$ & \textcolor{gris}{2} & \textcolor{gris}{4} & \textcolor{gris}{4} & \textcolor{gris}{4} & \textcolor{gris}{4} & \textcolor{gris}{4} & \textcolor{gris}{4} & \textcolor{gris}{4} & \textcolor{gris}{4} & \textbf{7} & 4 & 4 \\
$\mathbf{10}$ & \textcolor{gris}{4} & \textcolor{gris}{2} & \textcolor{gris}{4} & \textcolor{gris}{4} & \textcolor{gris}{4} & \textcolor{gris}{4} & \textcolor{gris}{4} & \textcolor{gris}{4} & \textcolor{gris}{4} & \textcolor{gris}{4} & \textbf{7} & 4 \\
$\mathbf{11}$ & \textcolor{gris}{4} & \textcolor{gris}{4} & \textcolor{gris}{2} & \textcolor{gris}{4} & \textcolor{gris}{4} & \textcolor{gris}{4} & \textcolor{gris}{4} & \textcolor{gris}{4} & \textcolor{gris}{4} & \textcolor{gris}{4} & \textcolor{gris}{4} & \textbf{7} \\
\end{tabular}

\end{center}

La diagonal vale $w(b_i)=7$ en los doce casos, que es impar, y todas las
entradas de fuera de la diagonal son pares. Tomando paridad queda
$(B^2)_{ii}=1$ y $(B^2)_{ij}=0$ para $i\neq j$, es decir $B^2=I_{12}$.

\subsection{Dimensión, tasa y overhead}

\begin{align*}
  n &= 24, \qquad k = 12, \qquad n-k = 12 \quad\text{(bits de paridad)}\\
  R &= \frac{k}{n} = \frac{12}{24} = \frac{1}{2} = 0{,}5 \\
  \text{overhead} &= \frac{n-k}{k} = \frac{12}{12} = 1 = \text{100\,\%}
  \qquad\left(\text{o bien } \frac{n-k}{n}=\text{50\,\%}\ \text{del bloque}\right)
\end{align*}

El código es $(24,12)$: entran 12 bits de mensaje y salen 24 de palabra
código. Hay $2^{12}=4096$ palabras código.

\subsection{Codificación a mano de $m=\hextt{A5C}$}

\[
  m = \hextt{A5C} = \texttt{1010\,0101\,1100}
  \quad\Longrightarrow\quad
  m_i = 1 \text{ para } i \in \{0,2,5,7,8,9\}
\]

La paridad es $p=m\cdot B$, el XOR de esas seis filas:

\begin{center}\small
% generado por gen_parte_a.py -- no editar a mano
\begin{tabular}{@{}r@{\;\;}l@{\;\;}l@{}}
 & $b_{0}$ & \texttt{1001\,1000\,1111} \\
$\oplus$ & $b_{2}$ & \texttt{0011\,0101\,0111} \\
\cmidrule(lr){2-3}
 & & \texttt{1010\,1101\,1000} \\
$\oplus$ & $b_{5}$ & \texttt{0111\,1100\,1100} \\
\cmidrule(lr){2-3}
 & & \texttt{1101\,0001\,0100} \\
$\oplus$ & $b_{7}$ & \texttt{0010\,1011\,1110} \\
\cmidrule(lr){2-3}
 & & \texttt{1111\,1010\,1010} \\
$\oplus$ & $b_{8}$ & \texttt{1000\,0111\,1011} \\
\cmidrule(lr){2-3}
 & & \texttt{0111\,1101\,0001} \\
$\oplus$ & $b_{9}$ & \texttt{1110\,0111\,0100} \\
\cmidrule(lr){2-3}
 & $p$ & \texttt{1001\,1010\,0101} \;$=$\; \texttt{9A5} \\
\end{tabular}

\end{center}

\[
  v = (\,m \mid p\,) = \texttt{1010\,0101\,1100}\;\texttt{1001\,1010\,0101}
    = \cwEj
\]

\subsection{Verificación $v\,H^{T}=0$}

Partiendo $v=(m\mid p)$ y usando otra vez $H^{T}=\left[\begin{smallmatrix}B^{T}\\ I\end{smallmatrix}\right]$:
\[
  v\,H^{T} \;=\; m\,B^{T} + p\,I \;=\; m\,B + p \;=\; p + p \;=\; 0 .
\]

Numéricamente, $m\cdot B = \parEj$ (la cuenta de arriba) y
$p = \parEj$, entonces
\[
  \texttt{1001\,1010\,0101} \oplus \texttt{1001\,1010\,0101}
  = \texttt{0000\,0000\,0000} = \hextt{000}. \qquad \checkmark
\]

\section{Ejercicio 2 --- Distancia mínima y capacidad de corrección}

Distribución de pesos del código (dato del enunciado):

\begin{center}\small
\begin{tabular}{@{}lccccc@{}}
\toprule
peso $w$ & 0 & 8 & 12 & 16 & 24 \\
\midrule
cantidad de palabras & 1 & 759 & 2576 & 759 & 1 \\
\bottomrule
\end{tabular}
\end{center}

El código es lineal, así que la distancia mínima es el peso mínimo no nulo:
\[
  d_{\min} = \min\{\,w(v) : v \in C,\; v \neq 0\,\} = 8 .
\]
(Chequeo: $1+759+2576+759+1 = 4096 = 2^{12}$, están todas las palabras.)

\textbf{Corrección.} Para corregir $t$ errores las esferas de radio $t$
alrededor de cada palabra código no se tienen que tocar, lo que pide
$d_{\min}\ge 2t+1$:
\[
  2t+1 \le 8 \;\Longrightarrow\; t \le \tfrac{7}{2} = 3{,}5
  \;\Longrightarrow\; t = \left\lfloor \frac{d_{\min}-1}{2} \right\rfloor
  = \left\lfloor \frac{7}{2} \right\rfloor = 3 .
\]

\textbf{Detección.} Usando el código solo para detectar, cualquier patrón de
error de peso $\le d_{\min}-1 = 7$ no puede llevar una palabra código a otra,
así que se detectan hasta \textbf{7} errores.

\textbf{Los dos a la vez.} Corrigiendo $t$ y detectando $l\ge t$ hace falta
$d_{\min}\ge t+l+1$:
\[
  l \le d_{\min} - t - 1 = 8 - 3 - 1 = 4 .
\]
O sea: corrige 3 y detecta 4. Esto se ve en el Ejercicio 3 con $r_3$, que
tiene un patrón de error de peso 4: el decodificador no lo corrige, pero lo
marca como no corregible en lugar de entregar una palabra equivocada.

\section{Ejercicio 3 --- Decodificación manual}

\textbf{Algoritmo (Parte C).} Con $s = B\,r_{[23:12]} \oplus r_{[11:0]}$ y
$q = s\cdot B$, se prueban los casos en orden y gana el primero que se cumple:

\begin{center}\small
\begin{tabular}{@{}clc@{}}
\toprule
caso & condición & patrón de error \\
\midrule
1 & $w(s)\le 3$ & $e=(0\mid s)$ \\
2 & $\exists\, i: w(s\oplus b_i)\le 2$ & $e=(u_i \mid s\oplus b_i)$ \\
3 & $w(q)\le 3$, con $q=s\cdot B$ & $e=(q\mid 0)$ \\
4 & $\exists\, i: w(q\oplus b_i)\le 2$ & $e=(q\oplus b_i \mid u_i)$ \\
5 & ninguna de las anteriores & \texttt{o\_uncorrectable} \\
\bottomrule
\end{tabular}
\end{center}

$u_i$ es el vector de 12 bits con un uno en la posición $i$. Los empates
en $i$ se resuelven por prioridad de índice (gana el $i$ más chico).


\subsection{$r_1 = \hextt{A5D9A6}$}

$r_{[23:12]} = \hextt{A5D} = \texttt{1010\,0101\,1101}$, con unos en
$i\in\{0,2,5,7,8,9,11\}$; $r_{[11:0]} = \hextt{9A6}$.

\begin{center}\small
\begin{minipage}[b]{0.46\textwidth}\centering
  \rAXor
\end{minipage}\hfill
\begin{minipage}[b]{0.50\textwidth}\centering
  \rASind
\end{minipage}
\end{center}

$w(s) = \rAPesoS > 3$, el caso 1 no aplica. Se buscan los doce candidatos
$s\oplus b_i$:

\begin{center}\small
\rATablaDos
\end{center}

Salta $i=11$ con peso 2, entonces es el \textbf{caso 2}:
\[
  e = (\,u_{11} \mid s\oplus b_{11}\,)
    = (\,\texttt{0000\,0000\,0001} \mid \texttt{0000\,0000\,0011}\,)
    = \hextt{001003}, \qquad w(e)=3 .
\]

\begin{center}\small
\rARes
\end{center}

Mensaje decodificado: $\hat{m} = \hextt{A5C}$, que es el mismo del
Ejercicio 1.

\subsection{$r_2 = \hextt{A5F9A4}$}

$r_{[23:12]} = \hextt{A5F} = \texttt{1010\,0101\,1111}$, con unos en
$i\in\{0,2,5,7,8,9,10,11\}$; $r_{[11:0]} = \hextt{9A4}$.

\begin{center}\small
\begin{minipage}[b]{0.46\textwidth}\centering
  \rBXor
\end{minipage}\hfill
\begin{minipage}[b]{0.50\textwidth}\centering
  \rBSind
\end{minipage}
\end{center}

$w(s) = \rBPesoS > 3$: no es el caso 1. Candidatos $s \oplus b_i$:

\begin{center}\small
\rBTablaDos
\end{center}

Ninguno baja de peso 2, así que tampoco es el caso 2. Se pasa a
$q = s\cdot B$:

\begin{center}\small
\rBQ
\end{center}

$w(q) = \rBPesoQ > 3$: no es el caso 3. Candidatos $q \oplus b_i$:

\begin{center}\small
\rBTablaCuatro
\end{center}

Aparece $i=11$ con peso 2, entonces es el \textbf{caso 4}:
\[
  e = (\,q\oplus b_{11} \mid u_{11}\,)
    = (\,\texttt{0000\,0000\,0011} \mid \texttt{0000\,0000\,0001}\,)
    = \hextt{003001}, \qquad w(e)=3 .
\]

\begin{center}\small
\rBRes
\end{center}

Mensaje decodificado: $\hat{m} = \hextt{A5C}$, otra vez el del Ejercicio 1.

\subsection{$r_3 = \hextt{A5C9AA}$}

$r_{[23:12]} = \hextt{A5C}$, con unos en $i\in\{0,2,5,7,8,9\}$ --- el mismo
mensaje del Ejercicio 1 ---; $r_{[11:0]} = \hextt{9AA}$.

\begin{center}\small
\begin{minipage}[b]{0.46\textwidth}\centering
  \rCXor
\end{minipage}\hfill
\begin{minipage}[b]{0.50\textwidth}\centering
  \rCSind
\end{minipage}
\end{center}

$w(s) = \rCPesoS > 3$: no es el caso 1. Candidatos $s \oplus b_i$:

\begin{center}\small
\rCTablaDos
\end{center}

Ninguno con peso $\le 2$: tampoco es el caso 2. Se calcula $q = s\cdot B$:

\begin{center}\small
\rCQ
\end{center}

$w(q) = \rCPesoQ > 3$: no es el caso 3. Candidatos $q \oplus b_i$:

\begin{center}\small
\rCTablaCuatro
\end{center}

Ninguno con peso $\le 2$: tampoco es el caso 4. Se llega al \textbf{caso 5},
\texttt{o\_uncorrectable} en alto, y $\hat{m}$ y $e$ no son válidas.

Comparando con la palabra del Ejercicio 1, $v=\cwEj$:
\[
  r_3 \oplus v = \hextt{A5C9AA} \oplus \hextt{A5C9A5}
  = \hextt{00000F} = \texttt{0000\,0000\,0000\,0000\,0000\,1111},
  \qquad w = 4 .
\]

El error real tiene peso 4, uno más que $t=3$. Es justo el caso del
Ejercicio 2: el código lo detecta pero no lo puede corregir.

\end{document}
